\scanpage{145}
Considérons deux éléments
\[
(16)\qquad \sigma^\nu u_\cdot=\sigma^\nu.(u_{-1},u_1),\quad\sigma^{\nu'}.u'_\cdot=\sigma^{\nu'}.(u'_{-1},u'_1)
\]
de $D_{\mathbb R\times\mathbb R}$, on voit de suite qu'ils ont $\overline{\mathrm m}$ effet sur $\widetilde X$ \deleted{\unsure{et}} (et $\overline{\mathrm m}$ seulement sur $\partial\widetilde X=\widetilde{\partial X}$) ssi ils sont égaux: $D_{\mathbb R\times\mathbb R}$ opère \uline{fidèlement} sur $\widetilde X$, et $\overline{\mathrm m}$ sur $\partial\widetilde X=\widetilde{\partial X}$.

Pour que les éléments (16) aient $\overline{\mathrm m}$ action sur $X$, il f. et s. qu'on ait
\[
(17)\qquad \nu=\nu'\quad\text{et}\quad u_s-u'_s\in\mathbb Z\quad\forall s\in\unsure{\mathbb I|\mathbb J},
\]
la première condition \inserted{$\nu=\nu'$} exprimant d'ailleurs que les deux opérations envisagées ont $\overline{\mathrm m}$ effet sur l'orientation de $X$. La condition $u_s-u'_s\in\mathbb Z\quad\forall s\in\unsure{\mathbb I|\mathbb J}$ s'explicite par
\[
(18)\qquad\left\{\begin{aligned}
\frac{u_1-u_{-1}}2&=\frac{u'_1-u'_{-1}}2\qquad\text{i.e.}\quad u_1-\unsure{u_{-1}^{4}|u_{-1}}=u'_1-u'_{-1}\\
\frac{u_1+\unsure{u_{+1}|u_{-1}}}2&\equiv\frac{u'_1+u'_{-1}}2\quad(\mathbb Z)\quad\deleted{\unsure{\text{i.e. }2u_{+1}}}
\end{aligned}\right.
\]
\deleted{Ecrivons} Prenons comme paramètres, pour repérer un élément \inserted{$\sigma^\nu.(u_{-1},u_{+1})$} de $D_{\mathbb R\times\mathbb R}$, non plus ($\nu$ et) $u_{-1}$ et $u_{+1}$, mais les coeff. $a=\frac{u_1-u_{-1}}2$ et $b=\frac{u_1+u_{-1}}2$ de (14), en posant
\[
(19)\qquad f_{a,b}\overset{\text{déf}}{=}(u_{-1},u_1)\qquad\text{avec}\quad\begin{aligned}u_{-1}&=-a+b\\u_1&=a+b\end{aligned}
\]
Alors on trouve
\[
(20)\qquad\deleted{\unsure{f}}f_{a,b}=f_{a',b'}\Longleftrightarrow a=a',\ b\equiv b'\ (\mathbb Z)
\]
\margin{145}
